% TOP_PROBLEM: {"id": 1267, "title": "Recamán's Sequence Completeness", "category_id": 1, "category": {"id": 1, "name": "number_theory", "display_name": "Number Theory", "description": "Properties of integers, prime numbers, Diophantine equations.", "slug": "number-theory", "order_index": 1, "created_at": "2026-07-31T15:26:25.670Z"}, "status": "open"}
% ATTEMPT_STATUS: unresolved
% Research attempt by GPT_6_Astra_Ultra, 1 October 2026.
% Canonical UnsolvedMath ID; original statements and sources preserved.
% FIRST_POSTED: 2026-10-01
% Imported from: unsolvedmath_additional_500.tex; UnsolvedMath ID 1267
% !TeX program = lualatex
% Compile twice: lualatex 1267.tex
% Self-contained: no external bibliography, images, or \input files.
% Numeric section labels and visible numbers are original UnsolvedMath IDs.
\documentclass[11pt,a4paper]{article}
\usepackage[margin=24mm,headheight=14pt]{geometry}
\usepackage{fontspec}
\setmainfont{Latin Modern Roman}
\setsansfont{Latin Modern Sans}
\setmonofont{Latin Modern Mono}
\usepackage{amsmath,amssymb,mathtools}
\usepackage{unicode-math}
\setmathfont{Latin Modern Math}
\usepackage{microtype}
\usepackage{enumitem,xurl,needspace}
\usepackage[dvipsnames]{xcolor}
\usepackage{hyperref}
\hypersetup{unicode=true,colorlinks=true,linkcolor=MidnightBlue,urlcolor=MidnightBlue,
 pdftitle={UnsolvedMath Problem 1267},pdfauthor={GPT 6 Astra Ultra},bookmarksnumbered=true}
\usepackage{bookmark,tocloft,titlesec,fancyhdr}
\setlength{\cftsecnumwidth}{4.2em}
\cftsetpnumwidth{2.5em}
\cftsetrmarg{4em}
\titleformat{\section}{\Large\bfseries\raggedright}{\thesection}{0.7em}{}
\pagestyle{fancy}\fancyhf{}
\fancyhead[L]{\small\sffamily GPT 6 Astra Ultra research attempt}
\fancyhead[R]{\small Original catalogue IDs}
\fancyfoot[C]{\thepage}
\setlength{\parindent}{0pt}
\setlength{\parskip}{5pt plus 1pt}
\setlength{\emergencystretch}{3em}
\setcounter{tocdepth}{1}\setcounter{secnumdepth}{1}
\setlist[itemize]{leftmargin=3em,itemsep=4pt,topsep=4pt,parsep=0pt}
\allowdisplaybreaks
\newcommand{\N}{\mathbb{N}}
\newcommand{\Z}{\mathbb{Z}}
\newcommand{\Q}{\mathbb{Q}}
\newcommand{\R}{\mathbb{R}}
\newcommand{\C}{\mathbb{C}}
\newcommand{\F}{\mathbb{F}}
\newcommand{\A}{\mathbb{A}}
\newcommand{\PP}{\mathbb{P}}
\newcommand{\E}{\mathbb{E}}
\newcommand{\eps}{\varepsilon}
\newcommand{\dd}{\,\mathrm d}
\newcommand{\sref}[2]{\hyperlink{source:#1:#2}{[S#2]}}
\newif\ifshowcatalogue
\showcataloguetrue
\newcommand{\cataloguescope}[1]{\ifshowcatalogue
 \paragraph{Statement recorded in the supplied source.}
 #1\fi}
\begin{document}
% UnsolvedMath ID 1267; source catalogue code NT-060

\setcounter{section}{1266}

\section{Completeness of Recamán's Sequence}\label{1267}

{\small\sffamily UnsolvedMath ID 1267 · Number Theory}

\subsection{Definitions and mathematical statement}

\paragraph{Shared notation.}Unless a section says otherwise, $\N=\{1,2,\ldots\}$,
$\Z$ is the ring of integers, $\Q,\R,\C$ are the rational, real, and complex
fields, $[N]=\{1,\ldots,\lfloor N\rfloor\}$, and $\log$ is the natural
logarithm. For real functions of a parameter tending to infinity,
$f\ll g$ and $f=O(g)$ mean $|f|\leq Cg$ eventually for some fixed $C$;
$f\gg g$ reverses this comparison; $f=o(g)$ means $f/g\to0$;
$f\sim g$ means $f/g\to1$; and $f\asymp g$ means both order bounds.
A subscript on $O$ or $\ll$ specifies allowed constant dependence.
The expression $n^{o(1)}$ in an upper bound means growth smaller than
$n^\epsilon$ for each fixed $\epsilon>0$. Local definitions govern any
exception, finite-field convention, or use of the same symbol for another
object. An asymptotic claim is not automatically an effective algorithm.



Set $a_0=0$. For $n\geq1$, take $a_n=a_{n-1}-n$ when that number is positive and has not appeared among $a_0,\ldots,a_{n-1}$; otherwise take $a_n=a_{n-1}+n$. Repetitions arising from the addition case are permitted. The target is $\{a_n:n\geq0\}=\Z_{\geq0}$, not monotonicity of this highly nonmonotone sequence. \sref{1267}{1} \sref{1267}{2}

\paragraph{Statement recorded in the supplied source.}
Does every nonnegative integer appear in Recamán's sequence? \sref{1267}{1}

\subsection{Short English statement}

Does the greedy subtract-if-new, otherwise-add rule eventually visit every nonnegative integer?

\subsection{Research attempt}

\paragraph{Scope and source check.}
The subtraction candidate must be positive and unused; the addition
case has no unused-value test. In particular repeated addition outputs
must not be rejected. The community-maintained sequence record [E1]
was inspected on 1 October 2026 for the convention and initial terms.
The analysis below uses the exact supplied recurrence and distinguishes
unboundedness of the values from coverage of all values.

\paragraph{Three elementary invariants.}
Induction gives $a_n>0$ for $n\ge1$, because an accepted subtraction
is positive and an addition is positive. Also
\[
 |a_n-a_{n-1}|=n,\qquad
 0\le a_n\le\frac{n(n+1)}2.
\]
The upper bound follows by always taking the larger possible update,
starting at zero. Modulo two, addition and subtraction agree, giving
\[
                  a_n\equiv\frac{n(n+1)}2\pmod2.
\]
Thus the value parity has period four in the index. This restricts
possible hitting times of a fixed value but leaves infinitely many
admissible indices for either parity.

The sequence is unbounded. If all values were at most $M$, every
consecutive difference would be at most $M$ in absolute value, contrary
to $|a_n-a_{n-1}|=n$ for $n>M$. This proof is unconditional and does
not use the novelty test. It does not establish that any particular
small missing value will ever be visited.

\paragraph{An exact late first-hit condition.}
Fix $m>0$ and suppose it has not appeared by time $n-1$. A subtraction
at time $n$ reaches it exactly when
\[
                            a_{n-1}=n+m.
\]
If this equality holds, the subtraction is positive and unused by
hypothesis, so it is necessarily taken. For $n\ge\max\{m,2\}$, an
addition cannot first reach $m$, since $a_{n-1}>0$ and hence
$a_{n-1}+n>m$. The initial exception is $a_1=1$. Therefore after
finitely many early times the
question of ever reaching $m$ is exactly a moving-line intersection
problem: must $a_{n-1}$ eventually equal $n+m$ while $m$ is missing?

In particular if a value $m$ is never reached, the graph of the sequence
avoids that line at every candidate late time. Unboundedness gives no
such intersection automatically. Even a sequence with infinitely many
values above and below a line can jump over its integer heights, since
the jumps here grow with $n$ rather than being unit steps.

\paragraph{What the least-missing-value argument actually says.}
If $m$ is the least positive value absent from a finite prefix, then
all positive subtraction candidates below $m$ have already been used.
At a time when such a candidate occurs, the rule therefore takes the
addition instead. This creates upward jumps near an already-filled
region, but supplies no monotone decrease of the distance to the
boundary $m$. The next step size changes, and the used set remembers
all earlier visits. A proof that ``the process must fill the next hole''
needs an invariant controlling that changing distance; the recurrence
alone has not supplied it.

One can track the moving coordinate $b_n=a_n-n$. The two updates are
\[
 b_n=b_{n-1}-(n+1)\quad\hbox{after subtraction},\qquad
 b_n=b_{n-1}+(n-1)\quad\hbox{after addition}.
\]
A late first hit of $m$ is equivalent to $b_{n-1}=m+1$. These
updates still have unbounded jump sizes and history-dependent signs.
They do not form a fixed finite-state walk to which recurrence in a
finite graph can be applied.

\paragraph{Repetitions are part of the actual sequence.}
The recurrence gives $a_{20}=42$ and $a_{24}=42$. At time $24$ the
subtraction from $a_{23}=18$ would be negative, so the addition returns
to forty-two despite its earlier appearance. A program that disallows
this repeated output computes a different sequence and cannot test the
catalogue target. The seen set is used only to decide whether a positive
subtraction is legal; it is updated after either type of step.

\paragraph{Executed computation.}
The script performs the exact recurrence through $n=100000$, retaining
the seen set and verifying the parity identity at every step. It
obtains $74253$ distinct values including zero, $25748$ repeated steps,
and maximum value $592216$. The least missing nonnegative integer in
that finite prefix is $61$. This last statement concerns the prefix,
not the infinite sequence: no permanence of the missing value is
claimed. The first twenty-six terms and these statistics are stored
in \texttt{results\_second.json}; the reproducer is
\texttt{checks\_second.py} under
\path{_Local/Erdos_problems/UnsolvedMath/attempt_audit_2026_10_01/number_theory/}.

\paragraph{Remaining gap and outcome.}
The attempt proves unboundedness and an exact criterion for a late
first visit, but no mechanism forces all those moving-line
intersections. It also gives no certified permanently missing integer.
\textbf{Outcome: UNRESOLVED.} Long growth, many distinct values, and
a changing least missing value remain compatible with either outcome.

\subsection{Sources}

{\small

\begin{itemize}

\item[\hypertarget{source:1267:1}{S1}] ULAM.AI, \emph{UnsolvedMath}, supplied \texttt{problems.json}, record 1267 (NT-060), statement and background. The embedded literature-review date is 2026-08-17; that is the archive’s review date, not a fresh certification. Dataset landing page: \url{https://huggingface.co/datasets/ulamai/UnsolvedMath}.

\item[\hypertarget{source:1267:2}{S2}] OEIS Wiki, Recamán's sequence (A005132). \url{https://oeis.org/wiki/Recam%C3%A1n%27s_sequence}. Bibliographic information imported from the supplied record.


\item[E1] OEIS Foundation and contributors, A005132,
Recam\'an's sequence. \url{https://oeis.org/A005132}.
Sequence definition and initial-term record inspected 1 October 2026;
the deductions and finite simulation above are supplied independently.

\end{itemize}

}

\label{end:1267}
\end{document}
